% Generated draft: requires mathematical and source review.
\documentclass[11pt]{article}
\usepackage[margin=25mm]{geometry}
\usepackage{fontspec}
\setmainfont{Latin Modern Roman}
\usepackage{amsmath,amssymb,mathtools,unicode-math}
\setmathfont{Latin Modern Math}
\usepackage{enumitem,needspace,xurl,hyperref}
\setlength{\parindent}{0pt}
\setlength{\parskip}{5pt}
\setlength{\emergencystretch}{3em}
\newcommand{\N}{\mathbb{N}}
\newcommand{\Z}{\mathbb{Z}}
\newcommand{\Q}{\mathbb{Q}}
\newcommand{\R}{\mathbb{R}}
\newcommand{\C}{\mathbb{C}}
\newcommand{\F}{\mathbb{F}}
\newcommand{\A}{\mathbb{A}}
\newcommand{\PP}{\mathbb{P}}
\newcommand{\E}{\mathbb{E}}
\newcommand{\eps}{\varepsilon}
\newcommand{\dd}{\,\mathrm{d}}
\begin{document}
\section*{UnsolvedMath ID: 10300012}
\textbf{Source problem code:} catalogue entry
\subsection{Definitions and mathematical statement}
Let \(M\) be a hyperbolic \(3\)-manifold and \(\Lambda\subset M\) a full genuine lamination in the sense of the cited Question 6.2. A codimension-one lamination is locally a product of a surface disk with a closed transversal set. The supplied record describes the relevant fullness/genuineness hypothesis by requiring a complementary region \(U\) that is an ideal-polygon bundle over a circle. Thus there is a bundle \(p:U\to S^1\) with ideal-polygon fiber; its core circle is an embedded circle \(K\subset U\) representing the circle direction, in the sense intended by the source.

For every such \((M,\Lambda,U)\), the question has two parts:
\begin{enumerate}
\item Is \(K\) isotopic in \(M\) to a closed geodesic for the hyperbolic metric? With the usual knot-theoretic meaning of isotopic, this asks for an ambient isotopy \(F_t:M\to M\), \(t\in[0,1]\), with \(F_0\) the identity and \(F_1(K)\) an embedded closed geodesic.
\item Does this core circle have a noncoalescable insulator family?
\end{enumerate}
The second term is not defined in the supplied record and must be recovered from the source's insulator framework. No equivalence or implication between the two questions is supplied. Likewise, any essentiality, regularity, orientability, compactness, or completeness assumptions implicit in the source's lamination and hyperbolic-manifold conventions remain to be checked. The first question concerns isotopy, which must not be weakened to existence of a geodesic in a free homotopy class.
\subsection{Short English statement}
For a hyperbolic three-manifold with a full genuine lamination and an ideal-polygon-bundle complementary region, ask whether the region's core is isotopic to a geodesic and whether it has a noncoalescable insulator family. The source presents these as open questions; the supplied current-status discussion reports partial structural tools and no definitive verification of this particular claim.
\subsection{Sources}
\begin{itemize}
\item Local UnsolvedMath registry, permanent ID \(10300012\), exact source code \texttt{catalogue entry}, ``Sublaminations and superlaminations --- Question 6.2.'' No individual proposer or proposed year is supplied.
\item Danny Calegari, \emph{Problems in foliations and laminations of 3-manifolds} (2002), Question 6.2: \url{https://arxiv.org/abs/math/0209081}. The supplied background describes an author-source list of 67 numbered questions.
\item Source-TeX extraction, accessed 2026-07-29. Source status NEEDS\_REVIEW: presented as open, with a current primary-literature status audit still required. The summary mentions split branched surfaces, insulators, Calegari, and Tao Li, without further bibliographic references; it reports no definitive verification of this specific claim.
\item Dataset attribution: Ulam AI / UnsolvedMath Contributors, CC BY 4.0. Source rights NEEDS\_REVIEW: publicly accessible author/publisher copy, redistribution terms requiring review.
\end{itemize}
\end{document}
